\subsection{\acrlong{APR}}
\label{subsec:analyse_prelem}

L’\acrfull{APR} constitue la première étape de la sûreté de fonctionnement. Elle vise à identifier les situations pouvant
compromettre la sécurité du vol, l’intégrité du lanceur, ou la conformité aux exigences du \acrshort{CDC} \gls{Fusex} bi-étage.
Cette analyse globale est réalisée avant toute étude détaillée, afin de mettre en évidence les phases critiques du vol et les
sous-systèmes nécessitant un contrôle renforcé.\\

Dans le cadre du projet SP-02, plusieurs éléments introduisent des risques spécifiques : la séparation passive par aérofreinage,
l’allumage conditionnel du second étage en vol, la structure en peau porteuse des deux étages, ainsi que les mécanismes
d’ouverture des parachutes et l’avionique. L’\acrshort{APR} vise à caractériser les situations dangereuses que ces sous-systèmes
pourraient générer au niveau du comportement global du lanceur.\\

Cette première étape aboutit à l’identification des \acrfull{ERPs}(section~\ref{subsubsec:ERPs}), c’est-à-dire les états finaux
dangereux qu’il convient d’éviter. Ils servent ensuite de base pour l’analyse par phase
(section~\ref{subsubsec:analyse_par_phase}) et pour la définition des modes de défaillance présentés en
section~\ref{subsec:modes_deffaillance}.

\subsubsection{\acrlong{ERPs}}
\label{subsubsec:ERPs}

Les \acrfull{ERPs} correspondent aux situations dangereuses qu’il est nécessaire d’éviter pour garantir un vol sûr et conforme.
Ils constituent les \textbf{états finaux indésirables} que des chaînes de défaillances pourraient produire, et non les pannes ou
dysfonctionnements eux-mêmes.\\

Leur rôle est de fournir une structure claire à l’analyse de sûreté : chaque mode de défaillance identifié ultérieurement sera
évalué en fonction de sa capacité à conduire à l’un de ces événements, et la criticité associée sera examinée dans
l’\acrshort{AMDEC}. La règle ASF1 du \acrshort{CDC} Fusex bi-étage stipule que :

\begin{quote}
    \textit{L’événement redouté est l’écartement de la trajectoire réelle par rapport à la trajectoire simulée.} ~\cite[p.~21]{CDC_BiEtage}
\end{quote}

Les \acrshort{ERPs} définis pour le projet SP-02 sont présentés dans le tableau \ref{tab:ERP_SP02}.

\begin{table}[h!]
\centering
\begin{tabular}{|c|p{12.5cm}|}
\hline
\textbf{\acrshort{ERP}1} & 
\textbf{Déviation de trajectoire / sortie du gabarit.}\newline
Perte de contrôle ou trajectoire non conforme, provoquée par une instabilité, une perturbation aérodynamique, un défaut de
séparation, une rupture structurelle ou tout événement compromettant la dynamique nominale de vol. 
\\ \hline

\textbf{\acrshort{ERP}2} & 
\textbf{a - Allumage intempestif ou non autorisé du second étage en phase sol ou rampe.}\newline
Mise à feu du moteur du deuxième étage alors que la fusée n’est pas encore en vol ou dans un état non prévu pour supporter une
propulsion, représentant un risque majeur pour les personnes et les biens. \newline\newline
\textbf{b - Allumage intempestif ou non autorisé du second étage en vol.}\newline
Allumage en dehors des conditions requises (séparation non confirmée, attitude non acceptable, fenêtre temporelle inactive,
logique d’ordre erronée), lorsque la fusée est en vol et pouvant entraîner une perte de contrôle ou une violation du gabarit
de vol.
\\ \hline

\textbf{\acrshort{ERP}3} & 
\textbf{Non-déploiement ou déploiement dangereux du parachute (E1 ou E2).}\newline
Absence d’ouverture, ouverture tardive, partielle ou perturbée, conduisant à une chute à vitesse excessive, susceptible de
générer un impact dangereux.
\\ \hline

\textbf{\acrshort{ERP}4} & 
\textbf{Rupture structurelle majeure entraînant une perte de contrôle.}\newline
Dégradation critique de la structure (aileron, peau porteuse, jonctions, coiffe) pouvant provoquer une instabilité brutale,
une rotation excessive ou une dispersion hors gabarit.
\\ \hline
\end{tabular}
\caption{Liste des \acrfull{ERPs} du projet SP-02}
\label{tab:ERP_SP02}
\end{table}

\newpage

\subsubsection{Analyse des risques par phase de vol}
\label{subsubsec:analyse_par_phase}

L’analyse par phase de vol permet de déterminer, pour chaque étape du déroulement nominal, quels \acrshort{ERPs} sont
susceptibles de se produire et dans quelles conditions générales. Cette étude complète l’identification des \acrshort{ERP} en
les replaçant dans le contexte opérationnel du vol, sans encore entrer dans le détail des causes techniques qui seront
analysées en section \ref{subsec:modes_deffaillance} (Modes de défaillance).\\

Chaque phase du vol présente en effet des exigences particulières : stabilité initiale, transition correcte entre étages,
gestion de l’attitude, respect des logiques d’autorisation ou encore déploiement de la récupération. L’analyse suivante vise à
identifier les situations dans lesquelles une déviation de trajectoire, un allumage inapproprié, une rupture structurelle ou un
défaut de récupération pourraient survenir.

\begin{table}[h!]
\centering
\begin{tabular}{|c|p{1.3cm}|p{12cm}|}
\hline
\textbf{Phase} &
\textbf{\acrshort{ERPs}} &
\textbf{Caractéristiques et points d’attention} \\ \hline

\textbf{0} &
ERP2a &
Phase sol critique vis-à-vis de la sécurité pyrotechnique. Un allumage du second étage au sol ou en rampe constitue un risque
immédiat pour les opérateurs. Vérification de l’état des sécurités (MAF, SECR, SPYTC) et de l’intégrité avio/mécanique
indispensable. 
\\ \hline

\textbf{1} &
ERP1, ERP4 &
La stabilité aérodynamique doit être garantie durant l’accélération initiale et notamment directement en sortie de rampe. Une
rupture structurelle ou une instabilité importante peut entraîner une déviation de trajectoire. Les sollicitations mécaniques ne
doivent pas compromettre la future séparation ni l’intégrité des organes critiques. 
\\ \hline

\textbf{2} &
ERP1, ERP4 &
La séparation passive doit être franche et symétrique. Une séparation partielle ou perturbée peut conduire à une rotation
excessive ou une trajectoire non conforme. Une perte structurelle locale lors de cette phase peut également provoquer une
instabilité brutale. 
\\ \hline

\textbf{3} &
ERP1, ERP2b, ERP4 &
Phase la plus sensible pour un bi-étage. Un allumage intempestif (en dehors des conditions d’attitude, de séparation ou de
fenêtre temporelle) peut entraîner une perte de contrôle. Une rupture structurelle ou une instabilité à l’allumage peut
également provoquer une sortie du gabarit. 
\\ \hline

\textbf{4} &
ERP3 &
Le bon déclenchement de la trappe et le déploiement du parachute conditionnent la vitesse de descente. Un défaut de récupération
(déploiement tardif ou partiel) conduit à un risque d’impact violent. 
\\ \hline

\textbf{5} &
ERP3 &
L’atterrissage doit se faire dans le gabarit prévu. En cas de non-déploiement ou de déploiement perturbé, l’impact peut être 
dangereux pour les personnes et les biens. 
\\ \hline

\end{tabular}
\caption{Analyse des \acrshort{ERPs} par phase du vol de SP-02}
\label{tab:ERP_par_phase}
\end{table}

Cette analyse montre que les phases critiques sont principalement la séparation inter-étages (phase~2) et l’allumage du second
étage (phase~3), où plusieurs \acrshort{ERPs} peuvent se matérialiser simultanément en fonction des conditions de vol. Les
risques liés à la récupération (phases~4 et~5) constituent également un enjeu majeur pour la sûreté au sol. Ces éléments servent
de base à l’identification des modes de défaillance en section \ref{subsec:modes_deffaillance}.

\newpage
